\documentclass[11pt,a4paper]{article}
\usepackage[utf8]{inputenc}
\usepackage[T1]{fontenc}
\usepackage[french,english]{babel}
\usepackage{amsmath,amssymb,amsthm}
\usepackage{amsfonts}
\usepackage{mathtools}
\usepackage{geometry}
\usepackage{hyperref}
\usepackage{enumitem}
\usepackage{booktabs}
\usepackage{array}
\usepackage{mathrsfs}
\usepackage{xcolor}
\usepackage{microtype}
\usepackage{tikz-cd}

\geometry{a4paper, top=2.5cm, bottom=2.5cm, left=2.5cm, right=2.5cm}

\theoremstyle{plain}
\newtheorem{theorem}{Th\'eor\`eme}
\newtheorem{proposition}[theorem]{Proposition}

\theoremstyle{definition}
\newtheorem{remark}[theorem]{Remarque}
\newtheorem{example}[theorem]{Exemple}

\DeclareMathOperator{\Sw}{Sw}
\DeclareMathOperator{\dimtot}{dim.tot}
\DeclareMathOperator{\rg}{rg}
\DeclareMathOperator{\Spec}{Spec}
\DeclareMathOperator{\Gal}{Gal}
\DeclareMathOperator{\codim}{codim}

\begin{document}

% ============================================
% FIRST LETTER: 28/10/76 to Luc Illusie
% ============================================
\selectlanguage{french}

\begin{flushright}
28/10/76
\end{flushright}

\noindent Cher Luc,

\medskip

Voici une semi-continuit\'e que j'avais conjectur\'ee dans le temps sur les conducteurs de Swan.

\begin{theorem}
Soit $f: X \to S$ un morphisme lisse s\'epar\'e de dim relative $1$, $F \subset X$ fini sur $S$ et $\mathcal{F}$ un faisceau localement constant sur $U = X - F$, prolong\'e par $0$ sur $X$.

\medskip
(Pour ce qui nous int\'eresse, $\mathcal{F}$ peut \^etre simplement un faisceau de $\mathbb{Z}/\ell$-modules). On suppose $\mathcal{F}$ de rang constant $r$.

\medskip
On suppose $F \to S$ universellement ouvert.

\medskip
Pour $s \in S$, soient $\bar{s}$ un point g\'eom\'etrique au-dessus de $s$, et
\[
\varphi(s) = \sum_{f \in F_{\bar{s}}} \bigl(\Sw_f(\mathcal{F} \text{ sur } X_{\bar{s}}) + r\bigr).
\]

\medskip
Alors

\textup{(i)} la fonction $\varphi(s)$ est semi-continue inf\'erieurement, et constructible;

\textup{(ii)} si elle est constante, $(X, \mathcal{F})$ est universellement localement acyclique.
\end{theorem}

\medskip

\noindent \textit{Preuve.} Des arguments bien connus donnent que $\varphi$ est constructible. Ceci ram\`ene la preuve de la semi-continuit\'e au cas o\`u $S$ est un trait. Le r\'esultat suivant ram\`ene de m\^eme la preuve de l'acyclicit\'e locale au cas o\`u $S$ est un trait:

\begin{proposition}
Soient $S' \to S$ un morphisme local de sch\'emas strictement locaux, $\bar{\eta}'$ un point g\'eom\'etrique de $S'$, et $\bar{\eta}$ son image dans $S$. Soient $s'$ et $s$ les points ferm\'es. Soit $X$ sur $S$, et $X'$ sur $S'$ obtenu par changement de base.
\[
\begin{tikzcd}
X_{\bar{\eta}'} \arrow[r, "\bar{g}"] \arrow[d] & X \arrow[d] & X_{\bar{\eta}} \arrow[l] \arrow[d] \\
\bar{\eta}' \arrow[r, hook] & S & \bar{\eta} \arrow[l, hook]
\end{tikzcd}
\]
et de m\^eme avec $s'$.
\end{proposition}

Si $\mathcal{F}$ est un faisceau sur $X$, et que le lieu de non lissit\'e de $(X, \mathcal{F})/S$ coupe $X_{\bar{\eta}}$ en $x$ en un \underline{point isol\'e}.

Alors, pr\`es de $x$, $i'^* R\bar{j}'_* \bar{j}'^* \mathcal{F}'$ est l'image inverse de $i^* R\bar{j}_* \bar{j}^* \mathcal{F}$ (et est constructible).

\medskip

\noindent \textit{Preuve:} comme SGA~7~XIII~2.4.2 (ingr\'edients: r\'eduction au cas propre, Th d'acyclicit\'e locale pour affirmer que ``$R\phi$'' est \`a support dans $x$ au voisinage de $x$, et invariance de la cohomologie par changement de corps alg.\ cl.\ et th.\ de ch.\ de base propre pour relier ``cycles \'evanescents'' et col.\ de la fibre g\'en\'erale).

\medskip

Supposons donc $S$ un trait strictement local; les assertions (i) (ii) r\'esultent du r\'esultat suivant:

\begin{proposition}
Soit $f: X \to S$ lisse s\'epar\'e, $F \subset X$ fini plat sur $S$, et local, de point ferm\'e $x$, $\mathcal{F}$ localement constant de rang $r$ sur $X - F$, prolong\'e par $0$. On a
\[
\Sw(\mathcal{F} \text{ sur } X_{\bar{\eta}}, \text{ en } x) + r = \sum_{f \in F_{\bar{\eta}}} \bigl(\Sw(\mathcal{F} \text{ sur } X_{\bar{\eta}}, \text{ en } f) + r\bigr) - \dim R\Phi_x^1(\mathcal{F}).
\]
\end{proposition}

\medskip

\noindent \textbf{NB:} a) L'hypoth\`ese de platitude sur $F$ assure que $R\Phi_x^i(\mathcal{F}) = 0$ pour $i \neq 1$; c'est pour ceci qu'on a dans le Th demand\'e $F/S$ univ.\ ouvert.

b) pour prouver (i) + (ii), il suffit de prouver ce r\'esultat, avec $F$ connexe: pour $F = \coprod F_i$ en effet, on a
\[
\begin{cases}
\text{(i) pour les } F_i \Rightarrow \text{(i)}, \\
\text{(b) constance de } \varphi(s) \underset{\text{(ii)}}{\Rightarrow} \text{constance pour les } F_i, \ldots
\end{cases}
\]

c) Le cas qui m'int\'eresse le plus est celui o\`u $F$ est une section de $X/S$. Dans ce cas, on trouve que $\Sw(\mathcal{F} \text{ sur } X_{\bar{\eta}}, \text{ en } g(\bar{\eta}))$ est semi-continue inf\'erieurement. Toutefois, on peut avoir \'egalit\'e constante de $\varphi$, et ne pas \^etre dans ce cas: deux points quadratiques ordinaires, en car.\ $0$, ($\Sw = 0$) peuvent confluer en un point quadratique non d\'eg\'en\'er\'e, en car.\ $2$ ($\Sw = 1$):

\medskip

De fa\c{c}on moins elliptique, consid\'erer sur $\mathbb{A}^1/\Spec(\mathbb{Z}_2)$ le faisceau d\'efini par le rev\^etement double $T^2 + bT + c = 0$, avec $a$, mod $2$, $b$ s'annule simplement en $0$.

\medskip

\noindent \textit{Lemme:} Gr\^ace au Th de r\'eduction semi-stable, on se ram\`ene \`a supposer que $X$ se plonge dans $\bar{X}$ sur $S$.

\medskip

\noindent \textit{Preuve:} Il nous est loisible de ramifier $S$. Ceci nous ram\`ene \`a supposer que $X$ se plonge dans $\bar{X}$ sur $S$, propre, et lisse sauf pour un nombre fini de points quadratiques ordinaires en r\'eduction (Th de r\'eduction semi-stable).

Le faisceau $\mathcal{F}$ donne lieu sur $\bar{X}$ \`a un faisceau loc.\ c$^{\text{st}}$ sur $\bar{X} - F$ --- partie de la fibre sp\'eciale --- partie plate sur $S$, ferm\'es ne contenant pas $x$.

\medskip

Nous allons maintenant passer de $\bar{X}$ \`a $X'/\bar{X}$; on veut

a) $X'$ est une courbe comme $\bar{X}$, et \'etal\'e au-dessus de $x$;

b) sur $X'$ --- (image r\'eciproque de $F$), $\mathcal{F}$ se prolonge en un faisceau loc.\ c$^{\text{st}}$.

\medskip

On le fait par \'etapes:

\medskip

$X'/S$:

a) trouver $X'/\bar{X}$ \'etal\'e au-dessus de $x$, propre sur $\bar{X}$, au-dessus duquel $\mathcal{F}$ se prolonge en dehors de $F'$: ce $X'$ sera d\'efini comme un ``joint'', de sorte qu'il suffit de trouver de tels $X'$ qui, localement sur $x$, mangent ainsi la ramification de $\mathcal{F}$. Le point est que si $\sum a_i T^i = 0$ d\'efinit un rev\^etement de $U$, ramifi\'e le long de $V$, alors, si on perturbe les $a_i$ par des quantit\'es tr\`es nulles le long de $V$, on trouve un rev\^etement qui, le long de $V$, est localement isomorphe au pr\'ec\'edent.

b) r\'eappliquer le Th de r\'eduction semi-stable.

\medskip

\noindent \textit{situation \'etudi\'ee, r\'ep\'et\'ee}

\newpage

Compasons $\chi(X'_{\bar{\eta}}, \mathcal{F})$ et $\chi(X'_{\bar{s}}, \mathcal{F})$.

\textcircled{1} Par la formule d'E.P.: si $X'/\bar{X}$ est \`a $h$ feuillets:
\begin{align*}
\chi(X'_{\bar{\eta}}, \mathcal{F}) &= \chi(X'_{\bar{\eta}}) \cdot r - h\Bigl(\sum_{f \in F_{\bar{\eta}}} (\Sw(\mathcal{F} \text{ sur } X_{\bar{\eta}}, f) + r)\Bigr), \\
\chi(X'_{\bar{s}}, \mathcal{F}) &= \chi(X'_{\bar{s}}) \cdot r - h\bigl(\Sw(\mathcal{F} \text{ sur } X_{\bar{s}}, x) + r\bigr).
\end{align*}

$\chi(X'_{\bar{s}}) = \chi(X_{\bar{s}}) \quad \# p^e$ doubles en fibre sp\'eciale

\begin{multline*}
\chi(X'_{\bar{\eta}}, \mathcal{F}) - \chi(X'_{\bar{s}}, \mathcal{F}) = -(\# p^e \text{ doubles sur } X_{\bar{s}}) \cdot r \\
- h\Bigl(\sum_{f \in F_{\bar{\eta}}} (\Sw + r) - (\Sw + r \text{ sur } X_{\bar{s}})\Bigr).
\end{multline*}

\textcircled{2} Par la m\'ethode des cycles \'evanescents: en $x$, on a $\Phi_x = \Psi_x$, et seul $\Psi_x^1 \neq 0$. En les points doubles apparus en fibre sp\'eciale, on a seulement $\Phi^1$, de dim.\ $r$.

Au total
\[
\Delta\chi = -(\# p^e \text{ doubles sur } X_{\bar{s}}) \cdot r - h(\Psi_x^1 \text{ pour } \mathcal{F}/X).
\]

Comparant ces formules, on trouve le r\'esultat annonc\'e.

\medskip

\noindent \textbf{Variante:} Soit $X \hookleftarrow F$ avec $F$ fini sur $S$, $X$ lisse de rel.\ dim.\ $1$ sur $S$ en dehors de $F$, $\mathcal{F}$ faisceau loc.\ c$^{\text{st}}$ de rang $r$ sur $X - F$.
\[
\begin{tikzcd}
X \arrow[d] \\
S
\end{tikzcd}
\]

\medskip

Alors: pour $\mathcal{F}$ virtuel de dimension $0$, on a
\[
\sum_{\substack{f \in F_{\bar{s}} \\ b \text{ branche par } f \text{ de } X_{\bar{s}}}}
\mspace{-20mu}
\Sw(\mathcal{F} \text{ sur } b) \quad
= \sum_{\substack{f \in F_{\bar{\eta}} \\ b \text{ branche par } f \text{ de} \\ X_{\bar{\eta}}}}
\mspace{-20mu}
\Sw(\mathcal{F} \text{ sur } b) \quad
+ \sum_{\substack{x \in F_{\bar{s}} \\ x \in F^c \\ \ldots}} \dim \Psi_x^*(\mathcal{F}).
\]

\medskip

\begin{flushright}
La preuve est la m\^eme.

\medskip

Bien \`a toi,

P.~Deligne
\end{flushright}

\vspace{1cm}

\noindent \textit{[Notation:} \textit{On a mentionn\'e aussi une variante avec} $\dimtot$\textit{.]}


\newpage
% ============================================
% SECOND LETTER: 28/10/76 (Troisieme lettre)
% ============================================

\begin{flushright}
28/10/76
\end{flushright}

\noindent Cher Luc,

\begin{flushright}
Troisi\`eme lettre (j'\'etais en forme hier)
\end{flushright}

\begin{theorem}
Soit $X$ propre sur $k$ alg.\ clos, et $\mathcal{F}_1, \mathcal{F}_2$ deux \'el\'ements du groupe de Grothendieck des faisceaux constructibles de $\bar{\mathbb{Q}}_\ell$-vectoriels. Si $\mathcal{F}_1$ et $\mathcal{F}_2$ sont localement \'egaux, alors $\chi(\mathcal{F}_1) = \chi(\mathcal{F}_2)$.
\end{theorem}

On proc\`ede par r\'ecurrence sur la dimension de $X$. Soit
\[
\begin{tikzcd}
X' \arrow[r, "g"] \arrow[d, "f"'] & X \\
\mathbb{P}'
\end{tikzcd}
\]

pour $X' \to X$ birationnel, la r\'ecurrence montre qu'il suffit de consid\'erer $X'$, $\mathbb{P}'$ et les images inverses $\mathcal{F}'_1, \mathcal{F}'_2$.

\[
\chi(X, \mathcal{F}) = \chi(X', \mathcal{F}') + \chi(\text{ferm\'e } F \text{ o\`u } \varphi \text{ n'est pas iso}, \mathcal{F}) - \chi(\varphi^{-1}F, \mathcal{F}'))
\]
et on regarde la suite spectrale de Leray de $f$. On a
\[
\chi(X', \mathcal{F}') = \chi(\mathbb{P}') \cdot \chi(\text{fibre g\'en\'erique}, \mathcal{F}) - \sum_{t \in \mathbb{P}'} \dimtot_p R\Gamma(X_t, \phi^* j_! \mathcal{F}))
\]

$g$ compos\'e un $\bar{}$, cf SGA~7.

On va voir que le $2^e$ membre est terme \`a terme pareil pour $\mathcal{F}_1$ et $\mathcal{F}_2$. Pour le $1^e$ terme, cela r\'esulte de l'hyp.\ de r\'ecurrence. Reste \`a voir
\[
\dimtot R\Gamma(X_t, \phi^*),
\]
vore plus g\'en\'eralement
\[
\chi R\Gamma(X_t, \phi^*)
\]
dans la cat\'egorie des modules munis d'une action de l'inertie. L'hypoth\`ese assure que, dans la cat\'egorie des faisceaux de modules avec action de l'inertie $I$, ou plut\^ot son groupe de Grothendieck, $R\phi(\mathcal{F}_1)$ et $R\phi(\mathcal{F}_2)$ sont localement \'egaux. Reste \`a disposer de l'hypoth\`ese de r\'ecurrence pour les $\bar{\mathbb{Q}}_\ell[I]$-faisceaux. Il suffit de noter que
\begin{align*}
K(\bar{\mathbb{Q}}_\ell[I]\text{-faisceaux constr.\ sur } X) &= \\
K(\bar{\mathbb{Q}}_\ell\text{-faisceaux sur } X) &\otimes K(\bar{\mathbb{Q}}_\ell[I]\text{-modules})
\end{align*}
$I$-module de base les $\bar{\mathbb{Q}}_\ell$ repr.\ irr\'ed.\ de $I$.

\newpage

\noindent \textbf{Remarque:} L'ingr\'edient essentiel de cette preuve est le th\'eor\`eme de finitude pour les faisceaux de cycles \'evanescents de SGA~$4\frac{1}{2}$.

\bigskip

Par ailleurs, je crois maintenant comprendre $\chi$ pour les surfaces, lorsque, en codimension $1$, on ne fait pas appara\^itre d'extension r\'esiduelle ins\'eparable (ce qui s'applique donc \`a la conjecture dans ma lettre \`a Katz). Je n'ai pas tout de d\'emonstration compl\`ete.

Soit $X$ une surface --- suppos\'ee normale pour simplifier --- sur $k$ alg\'ebriquement clos, $D$ un diviseur de Weil sur $X$, $U = X - D$ et $\mathcal{F}$ un faisceau localement constant sur $U$, de rang $r$.

\medskip

Par localisation en un point g\'en\'erique $\delta$ de $D$, on obtient un trait $X_\delta$, de corps des fractions $K$. Le faisceau $\mathcal{F}$ correspond \`a une repr\'esentation fid\`ele de $\Gal(K'/K)$, $K'$ extension galoisienne finie de $K$. On fait l'hypoth\`ese

\medskip

$(*)$ Les corps r\'esiduels des points du normalis\'e de $X_\delta$ dans $K'$, qui sont au-dessus de $\delta$, sont s\'eparables sur $k(\delta)$.

\medskip

Si $\bar{\delta}$ est le spectre d'une cl\^oture s\'eparable de $\delta$, et que $X(\bar{\delta})$ est la localis\'ee stricte de $X$ en $\bar{\delta}$, avec les m\^emes notations, $(*)$ devient

$(*')$ Les corps r\'esiduels de $K'$ co\"incident avec celui de $K$.

\medskip

On peut, sous cette hypoth\`ese, d\'efinir le \underline{conducteur de Swan} de $\mathcal{F}$ en $\bar{\delta}$ --- not\'e $\Sw(\mathcal{F} \text{ en } \bar{\delta})$ --- et il existe un ouvert lisse $D^\circ$ de $D$, tel que chaque fois qu'une courbe lisse $Y$ sur $X$ coupe \underline{transversalement} $D^\circ$ en un point $a$, sp\'ecialisation d'un point g\'en\'erique $\delta$, on ait
\[
\Sw(\mathcal{F}|_Y \text{ en } a) = \Sw(\mathcal{F} \text{ en } \bar{\delta}).
\]

\begin{flushright}
\begin{minipage}{0.5\textwidth}
\{Tout ceci vaut pour tout diviseur $D$ lisse c.\`a d.\ $X$ lisse, en toute dimension;\}
\end{minipage}
\end{flushright}

\newpage

Si $y$ bouge, et continue \`a couper $D$ dans $D^\circ$ transversalement, on est dans une situation d'acyclicit\'e locale (lettre pr\'ec\'edente). On peut le voir de fa\c{c}on directe: le point est que $\mathcal{F}$ se trivialise sur $X' \xrightarrow{g} X$, et que pour $y$ lisse coupant $D$ dans $D^\circ$, et transversalement, $g^{-1}y$ est encore lisse.

\medskip

\noindent \textbf{R\'esultat local:} on prend $X$ strictement local, essentiellement de type fini sur $k$.

\begin{theorem}
Soit $f: X \to \mathbb{A}^1$, envoyant le point ferm\'e $x$ sur $0$. On suppose que $f$ est lisse en dehors de $x$, et que $f|_D$ est \'etal\'e en dehors de $x$. Soient $D_i$ les composantes irr\'eductibles de $D$, $j$ l'inclusion de $U$ dans $X$, et $\Sw_i$ le conducteur de Swan de $\mathcal{F}$ au point g\'en\'erique de $D_i$. Alors, notant $\Psi_x^*$ la fibre en $x$ des faisceaux de cycles \'evanescents $\Psi$, l'expression
\[
\dimtot \Psi_x^* \bigl(j_! \mathcal{F} - r j_! \underline{\mathbb{F}}_\ell + \sum \Sw_i \cdot \mathbb{F}_{\ell D_i}\bigr)
\]
est ind\'ependante de $f$.
\end{theorem}

\medskip

On le notera $\Sw_x(\mathcal{F})$.

\medskip

\noindent \textbf{Remarques:} (i) $\dimtot$ est la somme
\[
\sum (-1)^i (\dim \Psi^i + \Sw \Psi^i);
\]
l'expression sous $\Psi_x^*$ est un $\mathbb{F}_\ell$-faisceau virtuel. L'expression $\Psi_x^*$ ne d\'epend que de sa fibre g\'en\'erique; celle-ci, regard\'ee comme faisceau virtuel sur une courbe, a les propri\'et\'es suivantes:

\qquad a) rang nul au point g\'en\'erique

\qquad b) en tout point, $(\text{rang g\'en\'erique} - \text{rang de la fibre au point} + \Sw) = 0$.

\medskip

\noindent \textbf{R\'esultat global:} on prend $X$ propre.

\newpage

\begin{theorem}
Supposons $X$ propre, et soient $D_i$ les composantes irr\'eductibles de $D$. Soit $D_i^\circ = D_i \cap D^\circ$, et $\Sw_i$ comme pr\'ec\'edemment. On a
\[
\chi(X, j_! \mathcal{F}) = \chi_c(U) \cdot \rg \mathcal{F} - \sum_i \chi(D_i^\circ) \cdot \Sw_i - \sum_{x \in D - D^\circ} \Sw_x(\mathcal{F}).
\]
o\`u $\chi_c$ d'apr\`es Grothendieck.
\end{theorem}

\medskip

\noindent \textit{Esquisse de preuve:} on \'ecrit le Th2 sous la forme
\[
\chi(X, j_! \mathcal{F} - r j_! \underline{\mathbb{F}}_\ell + \sum \Sw_i \cdot (\mathbb{F}_\ell)_{D_i^\circ}) = \sum_{x \in D - D^\circ} \Sw_x(\mathcal{F}).
\]

Calculons $\chi$ \`a l'aide d'un pinceau de sections hyperplanes. Si les sections sont en g\'en\'eral lisses, transverses \`a $D^\circ$, et toutes g\'en\'eriquement lisses, la formule d'Euler--Poincar\'e donne la formule voulue, avec un $\Sw_x$ calcul\'e en terme du pinceau. Quand on varie le pinceau, deux probl\`emes, si on fixe $D^\circ$ comme avant:

a) les tangentes entre $D^\circ$ et les sections du pinceau

b) la variation de $f$ (cf Th1) li\'ee aux sections passant par les points de $D - D^\circ$.

\medskip

Pour a), on d\'emontre:

(2) si $x \in D_i^\circ$, alors $\Sw_x(\mathcal{F}) = \Sw_i$.

\medskip

On commence par prendre des pinceaux g\'en\'eraux de haut degr\'e, de degr\'e variable. Ceci fait varier le $\#$ de $p^e$ de tangences de $D^\circ$ avec le pinceau, et on trouve que (1) n'est tenable que si (2) est vrai pour une tangente g\'en\'erale. Ceci acquis, on v\'erifie Th1 par passage du global au local (cf lettre pr\'ec\'edente) puis on trouve le Th1.

\newpage
% ============================================
% THIRD LETTER: 3/7/79 to Fulton
% ============================================
\selectlanguage{english}

\begin{flushright}
3/7/79
\end{flushright}

\noindent Dear Fulton,

\medskip

Your proof of ``Zariski's theorem'' is beautiful! It seems to me that a paraphrase of your proof gives the result with the topological $\pi_1$, over $\mathbb{C}$. This letter is to make it clear to me; I suspect you already observed it.

Let $C \subset \mathbb{P}^2(\mathbb{C})$ be a nodal curve. Let $D$ be an irreducible component of $C$, and $\widetilde{D}$ = normalisation of $D$. Let $V(\widetilde{D})$ be a ``tubular neighbourhood'' of $\widetilde{D}$ in $\mathbb{P}^2(\mathbb{C})$:
\[
\widetilde{D} \hookrightarrow V(\widetilde{D}) \xrightarrow{\varphi} \mathbb{P}^2(\mathbb{C}).
\]

\medskip

\noindent \textbf{Aim:}
\[
\pi_1(V(\widetilde{D}) - \varphi^{-1}(C)) \longrightarrow \pi_1(\mathbb{P}^2 - C) \qquad \text{(1)}
\]

If (1) is granted, Abhyankar's arguments work.

To get (1), my method has been to take your proof, and to systematically look at what it was implying on $\pi_1$, by applying your statements to ramified coverings --- the removal ``algebraic''. This translates your key statement into:

\begin{theorem}
Let $Z$ be a smooth connected \underline{locally closed} subvariety of $(\mathbb{P}^d)^n$, with $\dim Z > d(n-1)$. The diagonal $\Delta$ has then a fundamental system of neighbourhood $V(\Delta)$, such that $Z \cap V(\Delta)$ is connected, and
\[
\pi_1(Z \cap V(\Delta)) \longrightarrow \pi_1(Z).
\]
\end{theorem}

\medskip

Th $\Rightarrow$ Aim: take $d = 2$, $n = 2$, $Z = (\mathbb{P}^2 - C) \times (D - \text{double points of } C \text{ on } D)$.

Put also $Z_1 = (\mathbb{P}^2 - C) \times \widetilde{D}$ (above $(\mathbb{P}^2)^2$);
\[
\begin{tikzcd}
V(\Delta) \cap Z \arrow[r] \arrow[d] & V(\Delta) \cap Z_1 \arrow[d] \\
Z \arrow[r] & Z_1
\end{tikzcd}
\qquad ; \text{ on } \pi_1: \quad 1 \longrightarrow 1 \text{ (true)}.
\]

and $V(\Delta) \cap Z_1 \approx V(\widetilde{D}) - \varphi^{-1}(C)$, so we have get more than aimed for:
\[
\pi_1(V(\widetilde{D}) - \varphi^{-1}(C)) \longrightarrow \pi_1(\mathbb{P}^2 - C) \times \pi_1(\widetilde{D}).
\]

\newpage

In this case, the $\pi_0$ statement is trivial, and only $\pi_1$ matters. Because of this, I will only try to prove part of \textcircled{Th}:
\[
\pi_1(\text{some component of } V(\Delta) \cap Z) \longrightarrow \pi_1(Z);
\]
the $\pi_0$ should however not be difficult to get (anyway, you have it).

\medskip

As in your proof, one compares
\[
\begin{tikzcd}
\mathbb{P}^{nd} \supset L^d \arrow[dr, "\varphi_1"'] & & (\mathbb{P}^d)^n \supset \Delta \arrow[dl, "\varphi_2"] \\
& W \text{ (= join of two compactifications of } \mathbb{A}^{nd})
\end{tikzcd}
\]
and, once an infinity hyperplane is chosen in $\mathbb{P}^d$, and $\mathbb{P}^{nd}$, with
\[
\mathbb{P}^{nd} - \infty = \mathbb{A}^{nd} = (\mathbb{A}^d)^n = (\mathbb{P}^d - \infty)^n,
\]
\[
\begin{cases}
\varphi_1 = \text{to blow up } (\infty\text{-hyperplane}) \cap \text{closure of factors } \mathbb{A}^d \text{ of } (\mathbb{A}^d)^n = \mathbb{A}^{nd} \text{ / disjoint from } L \\
\varphi_2 = \text{`` \quad '' } \infty\text{-hyperplane of } \Delta
\end{cases}
\]

Hence: for $V_2(\Delta)$ neighbourhood of $\Delta$ in $(\mathbb{P}^d)^n$, there is $V_1(L)$, neighbourhood of $L$ in $\mathbb{P}^{nd}$, with
\[
\varphi_1^{-1}(V_1(L^d)) \subset \varphi_2^{-1}(V_2(\Delta))
\]

We may replace $Z$ by $Z \cap \mathbb{A}^{nd}$ (because $\pi_1(Z \cap \mathbb{A}^{nd}) \to \pi_1(Z)$), and hence to prove \textcircled{Th}, it suffices to consider $\mathbb{P}^{nd} \supset L^d$, instead of $(\mathbb{P}^d)^n \supset \Delta$.

\begin{theorem}[$*$]
Let $Z \subset \mathbb{P}^N$ be a smooth connected locally closed subvariety, $L$ a linear subspace, and assume $\dim Z > \codim L$. Then, $L$ has a fundamental system of neighbourhood with $Z \cap V(L)$ connected and $\pi_1(Z \cap V(L)) \to \pi_1(Z)$.
\end{theorem}

As before, I will only care for $\pi_1$.

\newpage

A neighbourhood of $L$ in $\mathbb{P}^N$ contains all $L'$ close to $L$. Hence the reformulation:

\begin{theorem}[$**$]
Same assumption, with now $L$ general (= in a suitable Zariski open subset of the grassmannian). Then, $\pi_1(Z \cap L) \to \pi_1(Z)$.
\end{theorem}

For Th$^{**}$, one can proceed by induction, and be reduced to the case where $L$ is an hyperplane. I hope this case to be in the litterature; at least one can reduce this case to the one treated by L\^e and Hamm:
\textit{un th\'eor\`eme de Zariski de type Lefschetz}, Ann.\ Sci.\ ENS 6 3 (1973) p.~317--366, where they take $Z = (\mathbb{P}^N - \text{some hypersurface})$: to get the result for our $Z$, project it generically onto a linear subspace $\mathbb{P}^{\dim Z}$, and use the complement of the ramification locus as the $Z$ of L\^e: his result gives what I need for hyperplanes through the center of projection.

\medskip

In fact, much more than Th$^{**}$ should be true.

\medskip

\noindent \textbf{Conjecture} (perhaps well known to L\^e): One considers
\[
Z \text{ smooth, connected (not assumed proper)}, \qquad f: Z \longrightarrow \mathbb{P}^N, \qquad L \text{ linear subspace},
\]
and
\[
\pi_i(\ldots / V(L)) \longrightarrow \pi_i(Z)
\]
for \textit{tubular neighbourhood} $V(L)$ of $L$, of usual shape. This map should be bijective for
\[
i < \dim f(Z) - \codim L - \sup_k \bigl(2k - \codim \text{nf}(Z) \text{ of the locus where dim fibre} \geq k + \dim \text{generic fibre of } Z \to f(Z)\bigr)
\]
and surjective for $i = \dim f(Z) - \codim L$.

\medskip

Here are my reasons to hope for it.

a) a proof by Morse theory, if ok for a generic $L$, should work as well for any $L$, when $L$ is replaced by $V(L)$. Let us take for $L$ a general hyperplane.

\newpage

b) For the similar statement in cohomology, and generic $L$, one should prove a vanishing of the $R\nu^i$
\[
H_c^i(\mathbb{P}^N - L, Rf_! \mathbb{Z}),
\]
related by duality to the higher
\[
H^i(\mathbb{P}^N - L, Rf_! \mathbb{Z}),
\]
handled by Leray spectral sequence: the $f$ dim support $R^i f_! \mathbb{Z}$ is checked by $(R^i f_! \mathbb{Z})_s = 0$ if $i > 2 \dim f^{-1}(s)$, and $H^i(\mathbb{P}^N - L, \mathcal{F}) = 0$ for $i > \dim \text{support of } \mathcal{F}$.

\vspace{1cm}

\begin{flushright}
Best wishes,

P.~Deligne
\end{flushright}

\newpage
% ============================================
% FOURTH LETTER: 16 Nov 79 to Fulton
% ============================================

\begin{flushright}
16 Nov 79
\end{flushright}

\noindent Dear Fulton,

\medskip

I enjoyed very much your notes on connectivity \ldots

A few remarks:

\medskip

--- The proof of Cor 3 p24 can be expressed somewhat more algebraically, looking at the diagramm:
\[
\begin{tikzcd}
\pi_1(X_0 \times X_0 \cap V) \arrow[r] \arrow[d] & \pi_1(X_0 \times X_0) \arrow[d] \\
\pi_1(X \times X \cap V) \arrow[r] & \pi_1(X \times X) \\
\pi_1(X) \arrow[u, "\Delta s"]
\end{tikzcd}
\]

: the image of $\pi_1(X)$ in $\pi_1(X)^2$ covers $\Im(\pi_1(X_0) \to \pi_1(X))^2$ --- hence this $\pi_1(X_0) \to \pi_1(X)$ is trivial. The rest as in your paper.

\medskip

--- In the same vein, at the places where you assume ``$\widetilde{X}$ irreducible'', one can drop the assumption by weakening the conclusion to ``\ldots has an image containing $\Im(\pi_1(\text{normalisation of } X) \to \pi_1(X)) = \Im(\pi_1(\text{smooth locus of } X) \to \pi_1(X))$''. Of course, this $\Im(\pi_1 \to \pi_1)$ is $\pi_1(X)$ if and only if $\widetilde{X}$ is irreducible.

\medskip

--- misprints: p36, Cor 3 (B): assume $C' \neq \emptyset$.

p36, Cor 2, in char $p$: $V$ is ``tamely'' simply connected.

\medskip

--- p7 : p18 Cor (B): one could give a proof closer in spirit to [G--L] by proving a stronger theorem: if $X \to \mathbb{P}^1 - H$ is an etale covering, possibly with $n$ sheets, there are points of $H$ above which $\geq e$ sheets come together, if $e \leq d$, $n+1$.

Meaning: if one restrict the covering to a small ball around a variable $h \in H$, they are connected components of it that have $\geq e$ sheets.

\newpage

Proof: a) one can assume by induction that $d \geq n+1$, that over each hyperplane, it happens that $n$ sheets come together, and, by contradiction, that more never do.

This enables to extend $X \to \mathbb{P}^N - H$ to a ramified covering $\bar{X} \to \mathbb{P}^N$; above small balls, $\bar{X}$ is $H$ (finite covering). Let $R \subset \bar{X}$ be the locus where $n$ sheets come together: $R \xrightarrow{f} \mathbb{P}^N$ is finite-to-one, above small balls, $H$ (finite map), and $f(R) = Y$ is algebraic, and, $R$ is above $Y^\circ$ Zariski open in $Y$, $R$ is \'etale. One has to show that $\Delta \subset R \times_Y X$ is not open. Its openness would contradict
\[
\pi_1((\mathbb{P}^N - H) \times Y^\circ \cap V(\Delta)) \longrightarrow \pi_1((\mathbb{P}^N - H) \times Y^\circ).
\]

\medskip

--- Did Lazarsfeld prove the local analogue (= for a small punctured neighbourhood of $0$ in affine space) of your theorem with Hamm:

\medskip

\noindent \textbf{Suggestion}: take $Z \xrightarrow{f} \text{Ball} - \{0\}$, with $f$ ``algebraic looking''.
\[
Z \xrightarrow{\text{compactif.}} \bar{Z} \xrightarrow{\bar{f}} \text{Ball}.
\]
Then, assuming $Z$ irreducible, $L$ linear subspace (o.c.L.) with $\dim f(Z) - \codim L \geq 2$, one has, after a shrinking of the ball (depending on $L$):

(a) for $L$ general: $\pi_i(\bar{f}L) \xrightarrow{\sim} \pi_i Z$ for $i < n-1$, $\twoheadrightarrow$ for $i = n-1$

(b) for any $L$: $\pi_i(f^{-1}(V(L - 0))) \xrightarrow{\sim} \pi_i Z$ for \ldots (same).

\medskip

If this suggestion did hold, it would reprove the theorem. Indeed, write $\mathbb{P}^n = P(V)$ (set of directions in $V$), hence $(V - \{0\}) \to P(V)$.

Starting from $Z \to P(V) \times P(V)$, it pulls back to
\[
\begin{tikzcd}
Z \arrow[r] & P(V) \times P(V) \\
Z_1 \arrow[u] \arrow[r] & (V - \{0\}) \times (V - \{0\}) \arrow[u] \arrow[r, hookleftarrow] & V - \{0\}
\end{tikzcd}
\]

\newpage

while $\Delta \subset P(V) \times P(V)$ has one $J$: $\Delta_V \subset V \times V$, the graph of the identity map of $V - \{0\}$ --- and a linear subspace.

The suggestion, applied to $Z_1$ and $\Delta_V$, gives Fulton--Hansen.

\medskip

This leads to suggestions about the higher $\pi_i$. First a reprojective-isation of the above. Starting from
\[
Z \longrightarrow P(V) \times P(V),
\]
one introduces $Z_2 \subset P(V \times V)$: ($Z_1$ above / homotheties):
\[
P(V \times V) = \text{join }(P(V) \text{ and } P(V)) \supset \mathbb{G}_m\text{-bundle over } P(V \times V)
\]
($\mathbb{G}_m$-bundle = $P(V \times V) - P(V \times 0) - P(0 \times V)$), and take pull back of $Z$.

Then, Bertini applied to $Z_2 \subset P(V \times V)$ and to $P(\text{diagonal } V \subset V \times V)$ gives Fulton--Hansen for $Z \to P(V) \times P(V)$. Bertini for higher $\pi_i$.

\medskip

One suggests:

\medskip

\noindent \textbf{Conjecture/Suggestion}: Take $Z \xrightarrow{f} P(V) \times P(V)$, $f$ proper, $Z$ smooth irreducible. Then
\[
0 \to \pi_i f^{-1}\Delta \to \pi_i Z_2 \to Z \to \pi_i f^{-1}\Delta \to \pi_i Z \to 0
\]
\[
\pi_i f^{-1}\Delta \xrightarrow{\sim} \pi_i Z \qquad (i > 2)
\]
up to $i = \dim f(Z) - \codim \Delta$, the injectivity part being lost for the last $i$.

\footnote{I went overboard. This is for $f$ finite. In general, a correction like in the conjecture in my Bambah Colloq.}

\medskip

Of course, I want a variant with $f^{-1}V(\Delta)$ for non proper $Z$. The map $\pi_2 Z_2 \to Z$ is as follows: one $Z$, we have the line bundles $\mathcal{O}(1,0)$ and $\mathcal{O}(0,1)$. Look at the homotopy sequences of the corresponding $\mathbb{G}_m$-principal bundles: $\pi_2(Z) \xrightarrow{\partial} \pi_1(\mathbb{C}^*) = \mathbb{Z}$, and take the difference of the maps to $\mathbb{Z}$ for $\mathcal{O}(1,0)$ and $\mathcal{O}(0,1)$.

It vanishes on the image of $\pi_2(f^{-1}\Delta)$, because $\mathcal{O}(1,0)$ and $\mathcal{O}(0,1)$ become isomorphic on $f^{-1}(\Delta)$.

\medskip

--- I guess that Lazarsfeld observed that the proof of Cor p28 + his local theorem gives that for
\[
\begin{tikzcd}
X \arrow[r] & Y \arrow[r, hook] & \mathbb{P}^N
\end{tikzcd}
\qquad
\begin{minipage}[t]{0.4\textwidth}
$Y$ closed, smooth; $X/Y$ $d$-sheeted covering,
\end{minipage}
\]
if $2\codim Y + d \leq N$, then $X$ is simply connected.

What does one know about higher $\pi_i$ or $H_i$ --- possibly assuming $X$ smooth?

\vspace{1cm}

I am going to Moscow in a few days, for 2 weeks, and will take your ms with me, to leave it there. Could you send me a new one?

\vspace{1cm}

\begin{flushright}
Yours truly,

P.~Deligne
\end{flushright}

\end{document}
